\ifdefined\response\newpage\newpage\clearpage\pagebreak\else\fi
\label{sec:ass_d_note}
We shortly identify sufficient conditions that guarantee asymptotic optimality  (Assumption~\ref{ass:empirical_proc_mean}), in addition to standard conditions on the nuisance parameters $\what{\pi}_{-k,n},\what{\mu}^C_{-k,n}$ (\Cref{ass:overlap}, \Cref{ass:consistency}) and on outcomes (\Cref{ass:bounded_outcomes}). 
\begin{assumption}[Overlap]
\label{ass:overlap}
    For some $\eta>0$ and for all $k,n$, we have  
    \vspace{-0.2em}
    $$
    \eta\leq  \pi(X) \le 1-\eta, \quad \eta \leq \what \pi_{-k,n}(X) \leq 1-\eta \quad \text{a.s.}
    $$  
\end{assumption}
\vspace{-1.2em}
\begin{assumption}[Convergence rates of propensity and constrained outcome models]
\label{ass:consistency}
For all~$k \in \{1,\ldots,K\}$, 
both $\what\pi, \what\mu^C$ are consistent,\vspace{-0.3em}
    $$
    \|\what\pi_{-k,n} - \pi\|_{L_2(P)}=o_P(1), \quad  \|\what\mu_{-k,n}^C - \mu\|_{L_2(P)}=o_P(1)
    $$ 
$$\text{and also} \quad    \|\what\pi_{-k,n} - \pi\|_{L_2(P)} \cdot  \|\what\mu_{-k,n}^C - \mu\|_{L_2(P)}=o_P(n^{-\frac{1}{2}}).$$
\end{assumption}
\noindent As we discuss below, we can relax these assumptions
when guaranteeing double robustness (consistency under misspecified nuisance parameters).
As is typical, we assume that outcomes do not differ too much from their means, conditional on covariates. 
\begin{assumption}
[Outcomes close to conditional means] \label{ass:bounded_outcomes} For all $k,n$ and some $0<B<\infty$, 
\vspace{-0.3em}
$$\|A(Y-P[Y\mid X])\|_{L_2(P)} =\|A(Y-\mu(X))\|_{L_2(P)} \leq B. $$
\end{assumption}

We also require the following condition (\Cref{ass:empirical_proc_mean}) on the C-Learner outcome model. This condition is new to our setting, and warrants discussion. %

\begin{assumption}[Empirical process assumption]%
\label{ass:empirical_proc_mean}\ \vspace{-0.3em}
$$(P_{k,n}-P)(  \what\mu_{-k,n}^C(X)- \mu(X)) =o_P(n^{-1/2}).$$
\end{assumption}
\noindent In \cref{sec:other_methods}, we showed how versions of one-step estimation  and targeting can also be considered C-Learners. We show in \cref{sec:show_constant_shift} and \cref{sec:show_tmle} that cross-fitted versions of the self-normalized AIPW and TMLE for continuous and unbounded outcomes, under standard assumptions on $\what\pi_{-k,n}, \what \mu_{-k,n}$ (with $\what \mu_{-k,n}$ the usual outcome model learned using Equation~\eqref{eq:regular_outcome}) also satisfy Assumption~\ref{ass:empirical_proc_mean}. Thus, C-Learner results in Theorems~\ref{thm:asymp},~\ref{thm:dr} (to come) apply directly to these estimators. %
One future research direction is to understand the model classes and constrained optimization methods for which this assumption holds. 

\Cref{ass:empirical_proc_mean} is not implied by previous assumptions. Consider the following example: 

\begin{example}[Assumption \ref{ass:empirical_proc_mean} does not follow from Assumptions \ref{ass:overlap}, \ref{ass:consistency}, \ref{ass:bounded_outcomes}]\label{ex:ass_D}
Assume \Cref{ass:overlap}.
Let $\what\mu_{-k,n}$ be an \emph{unconstrained} solution to
$\what\mu_{-k,n}\in \argmin_{\tilde\mu\in\mathcal F} P_{-k,n}[A(Y-\tilde \mu(X))^2]$,
in contrast txo the constrained problem in \eqref{eq:constraint}. 
Then, define 
$\what\mu^C_{-k,n}(x)=\what\mu_{-k,n}(x)+\sum_{i=1}^{n/K} \mathbf{1} \left\{x=x_i\right\}$, with the sum of indicators over elements in the $P_{k,n}$ fold.
Then 
\begin{align}
\label{eq:01_counterex}
    P\left|\what\mu^C_{-k,n}(X)-\what\mu_{-k,n}(X)\right|=0\;\text{ and } \;
P_{k,n}\left(\what\mu^C_{-k,n}(X)-\what\mu_{-k,n}(X)\right)=1.
\end{align}
We make the following additional assumptions for this example. 
We also assume that $\|\what\mu_{-k,n}(X)-\mu(X)\|=o_P(1)$ (analogous to \Cref{ass:consistency}).
Let $\what\mu$ denote $\argmin_{\tilde \mu\in\mathcal F} P[A(Y-\tilde\mu(X))^2]$, the best fitted model over the entire population. For simplicity, assume that $Y$ is bounded to satisfy \Cref{ass:bounded_outcomes}, and that $\what\mu(X)$ is bounded as well. Then
\vspace{-0.7em}
\begin{align*}
(P_{k,n}-P)\left(\what\mu^C_{-k,n}(X)-\mu(X)\right)&=
(P_{k,n}-P)\left(\what\mu^C_{-k,n}(X)-\what\mu_{-k,n}(X)\right)
\\&+
(P_{k,n}-P)\left(\what\mu_{-k,n}(X)-\what\mu(X)\right)\\
&+(P_{k,n}-P)\left(\what\mu(X)-\mu(X)\right)  
\end{align*}
where the first term on the RHS is 1 from \eqref{eq:01_counterex}, the second is $o_P(1)$ by the cross-fitting lemma (\Cref{lem:cross_fitting} in \Cref{sec:show_for_aipw_tmle}), and the third is $o_P(1)$ by Chebyshev. 
Therefore, 
\vspace{-0.7em}
$$(P_{k,n}-P)\left(\what\mu^C_{-k,n}(X)-\mu(X)\right)=1+o_P(1)$$
so that 
Assumption~\ref{ass:empirical_proc_mean} does not hold for this example, while Assumptions \ref{ass:overlap}, \ref{ass:consistency}, \ref{ass:bounded_outcomes} do. 
\end{example}
\label{sec:ass_d_note-end}
\ifdefined\response\newpage\newpage\clearpage\pagebreak\else\fi
